\SourceLeaf{7}
\begin{center}
THEORIA ATTRACTIONIS CORPORUM\par
\bigskip
SPHAEROIDICORUM ELLIPTICORUM HOMOGENEORUM\par
\bigskip
METHODO NOVA TRACTATA.\par
\medskip
\rule{.16\textwidth}{.3pt}
\end{center}
\Article{1}
\par
Satis quidem constat, problema de attractione corporis sphaeroidici elliptici homogenei in punctum quodvis exacte determinanda ad quaestiones difficillimas astronomiae physicae referri, pluresque geometras, inde a \textsc{Newtoni} temporibus, acriter iteratisque vicibus illi incubuisse. Primo quidem, investigatione ad sphaeroidem per revolutionem semiellipsis circa alterutrum axem ortam restricta, ipse summus \textsc{Newton} attractionem quam patitur punctum in axi situm invenire docuit, simulque nexum inter attractiones, quas patiuntur puncta intra sphaeroidem in eadem diametro sita, assignavit (\textit{Princip. Lib.} I. \textit{Prop.} XCI). Dein sagax \textsc{Mac Laurin}, synthesi perelegante usus, attractionem punctorum in sphaeroidis superficie vel in prolongatione plani aequatoris positorum determinavit, quo pacto simul theoria attractionis punctorum intra sphaeroidem sitorum, quae per \textsc{Newtoni} theorema ad attractionem punctorum in superficie facile referebatur, complete absoluta erat (De caussa physica fluxus et refluxus maris, in \textit{Recueil des pièces qui ont remporté les prix de l’acad. roi. des sc. T.} IV; \textit{Treatise of fluxions B.} I. \textit{Ch.} 14). Quae \textsc{Mac Laurin} per synthesin enucleaverat, postea per analysin (cui antea huiusmodi quaestiones inaccessibiles visae erant) haud minus eleganter eruere docuit ill. \textsc{Lagrange}, atque sic viam ad ulteriores progressus patefecit (\textit{Nouv. Mém. de l’Acad. de Berlin} 1773). Scilicet adhuc desiderabatur attractio punctorum extra sphaeroidem neque vero in axis nec in aequatoris prolongatione sitorum enodanda,
